% Sección del estado del arte para el documento TT
% TT: Pipeline in silico usando GNN multimodal para la identificación de moléculas afines al receptor 5-HT1A (Ki, IC50, BBB, Tox)

% Como comienzo del capitulo del estado del arte se dará mediante un resuemen del capitulo donde explicamos lo hallazgos, lo que esté trabajo busca abordar y lo que se espera lograr con el desarrollo de este proyecto. Lo que se bysca resolver son los tres primeros puntos: 
% 1. ¿Por qué la predicción de afinidad, permeabilidad BBB y toxicidad es importante en el desarrollo de fármacos?
% 2. ¿Cuáles son los métodos actuales para predecir estas propiedades y cuáles son sus limitaciones?
%


\section{Estado del Arte}
\label{sec:state_of_the_art}

En la era contemporánea, la investigación en biología molecular y química computacional ha experimentado un crecimiento exponencial gracias a los avances en técnicas de modelado y simulación como lo es el \textit{big data }. La identificación de ligandos específicos para receptores biológicos, como el receptor 5-HT1A, es un área de gran interés debido a su relevancia en el desarrollo de fármacos y terapias para diversas enfermedades neurológicas y psiquiátricas. 


En los últimos años, el uso de modelos de aprendizaje profundo ha revolucionado la investigación en química computacional y biología molecular. En particular, las redes neuronales gráficas (GNN) han demostrado ser herramientas poderosas para predecir propiedades moleculares y actividades biológicas. Estas redes permiten representar moléculas como grafos, donde los átomos son nodos y los enlaces químicos son aristas, capturando así la estructura y las relaciones intramoleculares de manera efectiva.


\subsection{Predicción de afinidad para el receptor 5HT1A}
\label{ssec:arte_5ht1a}

El receptor 5-HT1A es un receptor acoplado a proteína G de clase~A que se expresa predominantemente en el sistema limbico y desempeña un papel crucial en la regulación de la serotonina en el sistema nervioso central. La afinidad de los ligandos hacia este receptor es un factor determinante en su eficacia terapéutica y selectividad.Ya que sólo el 5\% de la serotonina corporal actúa en el cerebro; el resto solo cumple funcione periféricas relacionadas con la motilidad intestinal y la agregación plaquetaria. No obstante, la fracción ceebral es farmacológicamente relevante, ya que el la activación o el bloqueo del receptor 5-HT1A modula la transmisión sináptica, la plasticidad neuronal, la neurogénesis y la neuroprotección ~\cite{czub_curated_2021}. Como consecuencia, la disfunción de este receptor se ha asociado con trastornos psiquiátricos como la depresión, la ansiedad y la esquizofrenia. Por lo tanto, la evidencia reciente sugiere que el receptor 5-HT1A puede constituir un objetivo terapéutico prometedor para el desarrollo de fármacos destinados a tratar estas afecciones, además de tener una correlación con el tratamiento para el cáncer colorrectal, de pulmón de células pequeñas y de próstata, mediante el bloqueo de la poliferación tumoral inducida por al serotonina ~\cite{czub_curated_2021}.


La naturaleza de ligando--receptor en los GPCRs (Proteínas acopladas a receptores G) hace que la selectividad sea un requisito fundammental para el desarrollo de fármacos que actúen sobre el receptor 5-HT1A. Asimismo el receptor 5-HT1A comparte similitudes estructurales con los recepetores 5-HT1B y 5-HT1D, lo que puede dar lugar a interacciones cruzadas no deseadas. Por lo tanto, los compuestos diseñados para este objetivo deben de descriminar con precisión entre estos receptores. Este requisito de selectividad, combinado con la necesidad de penetración del SNC, convierte la identificación de ligandos específicos para el receptor 5-HT1A en un desafío significativo en el diseño de fármacos. La predicción precisa de la afinidad de los ligandos es esencial para el diseño racional de fármacos, ya que permite identificar compuestos con alta probabilidad de éxito en etapas tempranas del desarrollo farmacológico. 

\subsubsection{Historia de los modelos QSAR para 5-HT1A}

Actualmente la modelización computacional de la actividad frente al receptor 5-HT1A tiene una historia de más de 3 décadas que, en su medidad, presentan limitaciones estructurales que este presente trabajo busca superar. La primera generación de los modelos QSAR (predominantes hasta mediados de la década de 2010) adoptó un enfoque de clasifiación binaria como lo es (activo/inactivo o agonista/antagonista) a paritr de descriptores fisicoquímicos y topológicos de las moléculas calculados con herramientas como MOE o Dragon. Estos modelos ofrecen valor heurístico para el cribado inicial, sin embargo,  carecían de la capacidad de capturar la complejidad de las interacciones ligando-receptor y no podían predecir con precisión la afinidad cuantitativa (Ki o IC50) de los compuestos, diferencia que es clínicamente decisiva ~\cite{czub_curated_2021}.

No obstante, la segunda generación intentó abordad la regresión sobre Ki o pKi, pero lo hizo con restricciones estructurales severas. Los conjuntos de entrenamiento estaban limitados a familias congenéricas como lo son arylpiperazas, tienopirimidona derivados u otras series sintéticamente homogéneas. Esta homgeniedad permite obtener coeficientes $R^2$ muy elevados (superando el valor de 0.90) y esto ocurría precisamente porque la variación química es pequeña y los modelos QSAR lineales capturan bien las tendencias dentro de una serie. Sin embargo, un modelo entrenado en estas familias congenéricas tiene dominio de aplicablidad virtualmente nulo sobre scaffolds distintos. lo que hace inapropiado el cribado virtual ~\cite{czub_curated_2021}.

\subsubsection{El trabajo de Czub et al. (2021): primera base de datos curada y diversa con pKi}

El avance más relevante en la modelización específica de 5-HT1A fue publicado por Czub, Pacławski, Szlęk y Mendyk en 2021 en \textit{Pharmaceutics}. La  contribución principal no es un modelo sino una \emph{base de datos}, que resuelve el cuello de botella que impedía la regresión diversa: la ausencia de un conjunto estructuralmente heterogéneo con pKi reportado para 5-HT1A.
 
La curación se integró con dos fuentes: ZINC y ChEMBL (descarga de noviembre
de 2020). La fusión produjo 11,649 compuestos iniciales; tras eliminar las moléculas duplicadas cuya diferencia de pKi entre bases superaba 0.1 unidades logarítmicas (criterio que elimina mediciones inconsistentes entre fuentes), el conjunto final contuvo 9,440 compuestos únicos, la colección diversa más grande de ligandos 5-HT1A con pKi reportado disponible públicamente hasta esa fecha~\cite{czub_curated_2021}. La distribución de pKi en el rango de 4.2 a 11.0 sigue una distribución gaussiana, lo que sugiere cobertura equilibrada del espacio de actividades.

Para la modelización, los autores emplearon \emph{AutoML} de la plataforma H$_2$O, que realiza búsqueda automática de hiperparámetros y selección de arquitecturas sobre un espacio de modelos que incluye redes neuronales profundas, \textit{gradient boosting machines} (GBM) y XGBoost. La representación molecular se basó en 1{,}826 descriptores Mordred 2D, reducidos a 216 variables mediante selección automática de características. El modelo ensamble final alcanzó $R^2 = 0.74$ y RMSE $= 0.54$ en validación cruzada de diez pliegues. La aplicación del método SHAP (\textit{Shapley Additive Explanations}) reveló que los descriptores de mayor importancia son de naturaleza electrostática y topológica, en particular descriptores de autocorrelación de Moreau--Broto ponderados por electronegatividad y parámetros geométricos del espacio molecular~\cite{czub_curated_2021}.


\subsubsection{El trabajo de Czub et al. (2022): Validación OECD y evaluación externa}

Después del trabajo de Czub en el 2021 publicada en \textit{Pharmaceutics}, desarrolló un trabajo en 2022 donde evaluó el modelo de AutoML contra los cinco principios de la OCDE para modelos QSAR de usos reglatorio como lo es: definición del punto final, algoritmo no ambiguo, dominio de aplicabilidad, validez estadística apropiada y, de ser posible, mecanismo de acción interpretable ~\cite{czub_automlbased_2022}

La contribución metodológica central es doble. En primer lugar, los autores emplearon el criterio de información de Akaike (AIC) (en lugar de métricas de rendimiento puro) para la selección del modelo final, lo que penaliza la complejidad excesiva y favorece la parquedad. En segundo lugar, condujeron 30 repeticiones independientes del proceso de entrenamiento (\textit{multi-start}) para verificar la reproducibilidad numérica del sistema AutoML: el análisis de varianza (ANOVA) de una vía confirmó que las diferencias entre corridas no son estadísticamente significativas, con RSD de 0.09–1.54\% en el conjunto curado y de 0.12–1.36\% en el conjunto externo GLASS. La validación externa sobre más de 700 moléculas de la base GLASS mostró capacidad predictiva real, no únicamente validación interna optimista. El modelo final es un ensamble de 17 modelos componentes (redes neuronales, GBM y XGBoost), disponible públicamente en GitHub ~\cite{czub_automlbased_2022}

Sin embargo, tanto el estudio del 2021 como el del 2022 emplean \emph{partición aleatoria} de los datos para la validación cruzada. Esta decisión contiene una consecuencia que lo autores no mencionan que es: en un conjunto de datos donde coexisten compuestos estructuramente similares, la particion aleatoria garantiza que el conjunto de prueba contendrá similitudes con el conjunto de entrenamiento dando así una inflación articial de las métricas de generalización. En el contexto de descubrimientos de fármacos donde el objetivo es predecr la actividasd de series \emph{nuevas} y no de compuestos similares a los ya conocidos, la partición mediante andamio o similitud representan e estandar metodológico adecuado. Dando así las métricas $R^2 = 0.74$ reportadas deben de interpretarse como un límite superior altruista. 

\subsection{Redes neuronales de grafos para predicción de afinidad medicamento--proteína}
\label{ssec:arte_gnn_dta}

\subsection{Del descriptor al Grafo}

Los modelos QSAR clásicos codifican la estructura molecular en un vector de longitud fija: descriptores fisicoquímicos calculados (logP, TPSA, descriptores de Moreau--Broto) o huellas de bits que marcan la presencia o ausencia de subestructuras predefinidas (ECFP4/Morgan). Esta representación tiene ventajas claras como lo es la eficiencia computacional, interpretabilidad, compatibilidad con métodos de ML estándar pero introduce una pérdida de información estructural que puede ser relevante para la predicción de actividad~\cite{jiang_could_2021}. 
 
Las GNNs adoptan una representación alternativa: la molécula como un grafo $G = (V, E)$ donde los nodos $V$ son los átomos y las aristas $E$ son los enlaces. Cada nodo y cada arista reciben vectores de características que codifican propiedades atómicas locales (número atómico, grado, carga formal, hibridización, aromaticidad, quiralidad) y propiedades del enlace (tipo, conjugación, estereoquímica). A diferencia del descriptor, esta representación preserva explícitamente la topología molecular y permite al modelo aprender representaciones contextuales de cada átomo en función de su entorno químico inmediato~\cite{jiang_could_2021}.
 
El mecanismo de aprendizaje en una GNN sigue un esquema: en cada capa, cada nodo agrega la información de sus vecinos, actualiza su representación propia mediante una función aprendible, y propaga la representación actualizada a la capa siguiente. Tras $k$ capas, la representación de un átomo incorpora información de su vecindad de radio $k$ en el grafo. Una operación de \emph{readout} final agrega las representaciones atómicas en un vector global que representa la molécula, y este vector alimenta capas densas de clasificación o regresión.


\subsubsection{SS-GNN: un grafo unificado para el complejo proteína--ligando}
 
Zhang, Jin, Liu y colaboradores propusieron en 2023 el modelo SS-GNN (\textit{Simple-Structured Graph Neural Network}), cuya contribución central es metodológica: en lugar de modelar proteína y ligando como grafos separados que luego se combinan —estrategia que requiere operaciones de fusión costosas y decisiones de diseño no triviales—, SS-GNN construye un \emph{único grafo de distancia umbralada} que codifica de forma conjunta los átomos del ligando y los residuos proteicos en el sitio de unión~\cite{zhang_ssgnn_2023}.
 
La representación se define a partir de un umbral de distancia entre átomos de ligando y proteína: sólo se incluyen interacciones dentro de un radio específico (evaluado entre 5 y 8~Å). Los enlaces covalentes dentro de la proteína son ignorados deliberadamente, lo que reduce de forma significativa el número de nodos y aristas sin pérdida relevante
de información para la predicción de afinidad. El módulo de extracción de características combina capas GIN (\textit{Graph Isomorphism Network}) con un perceptrón multicapa (MLP) que procesa independientemente las características de nodos (átomos) y de aristas (pares átomo-átomo), un mecanismo que los autores denominan DDA (\textit{Distance Discretization Aggregation}).
 
Con sólo 0.5 millones de parámetros (aproximadamente una décima parte de modelos equivalentes) y una velocidad de inferencia de 0.2 milisegundos por complejo, SS-GNN alcanzó una correlación de Pearson $R_p = 0.853$ en el conjunto \textit{core} de PDBbind v2016, superando a los métodos GNN previos en 5.2 puntos porcentuales~\cite{zhang_ssgnn_2023}. La eficiencia computacional es relevante en el contexto de cribado virtual: la posibilidad de procesar miles de complejos en segundos acerca este tipo de modelos a escenarios de uso real.
 
La limitación fundamental de SS-GNN es la dependencia de la estructura 3D del complejo proteína--ligando: el modelo requiere como entrada la pose de acoplamiento del ligando en el sitio de unión, que en cribado \textit{de novo} debe obtenerse mediante \emph{docking} molecular previo. Esto introduce un error acumulado (el error del \emph{docking} propagado al error de predicción de afinidad) que no está caracterizado en el trabajo original.



\subsubsection{Las GNNs como acelerador  del cribado a escala masiva}
 
Hosseini, Simini, Clyde y Ramanathan (2022), del Laboratorio Nacional de Argonne, abordan un problema diferente: dado que las librerías virtuales modernas contienen decenas o cientos de millones de moléculas, el \emph{docking} clásico (que tarda del orden de minutos por molécula en una CPU) es computacionalmente prohibitivo para el cribado exhaustivo~\cite{hosseini_deep_2022}.
 
Su propuesta, \textit{Deep Surrogate Docking} (DSD), reformula el problema: en lugar de sustituir el \emph{docking} por un modelo de ML,lo usa como oráculo para generar etiquetas de entrenamiento sobre unafracción pequeña de la librería, y luego entrena una GNN (arquitecturaFiLMv2) para predecir los \textit{docking scores} del resto. La GNNactúa así como un prefilador eficiente que prioriza los candidatos más prometedores para su evaluación con \emph{docking} real. El experimento de referencia utilizó 128 millones de moléculas de ZINC y el receptor de dopamina $D_4$ como diana. La arquitectura FiLMv2, que incorpora un mecanismo de atención por modulación de características de los vecinos (\textit{Feature-wise Linear Modulation}), logró una aceleración de $9.496\times$ en el tiempo de cribado con una tasa de error de recuperación inferior al 3\%: más del 97\% de los hits reales identificados por docking exhaustivo son recuperados por el pipeline DSD~\cite{hosseini_deep_2022}. Desde la perspectiva del presente proyecto, DSD ilustra un puntometodológico importante: las GNNs no deben evaluarse sólo como reemplazos de los métodos clásicos de predicción de afinidad, sino también como componentes de un pipeline más amplio donde la eficiencia y la escalabilidad son factores de diseño determinantes.
 
\subsubsection{La convergencia hacia modelos integrados}
 
Mao y Yan (2026), en una revisión de mini-alcance pero de amplio espectro temporal (enero 2007 -- abril 2026), documentan la trayectoria de convergencia entre arquitecturas de GNN, Transformers moleculares y modelos de difusión hacia pipelines integrados capaces de predecir simultáneamente propiedades de unión, ADMET y propiedades cristalográficas~\cite{mao_artificial_2026}. Los autores reportan que modelos de cardiotoxicidad basados en GNNs con huellas moleculares, entrenados sobre miles de compuestos de ChEMBL y PubChem, alcanzan AUC-ROC por encima de 0.90 en conjuntos de prueba independientes para la inhibición del canal hERG (predictor clave de arritmias cardiacas fatales), superando a los modelos QSAR clásicos en identificación de riesgos de seguridad tempranos~\cite{mao_artificial_2026}.
 
No obstante, Mao y Yan también señalan que los modelos actuales operan sobre características moleculares estáticas (peso molecular, área de superficie polar, lipofilia) ignorando la naturaleza dinámica y dependiente del contexto de las interacciones biológicas. La dosis, el crosstalk entre dianas, los gradientes de pH en el microentorno tumoral: ninguna de estas realidades biológicas está capturada en los modelos convencionales de ML molecular. Esta brecha epistemológica es, en última instancia, lo que limita la tasa de éxito en ensayos clínicos independientemente del rendimiento computacional~\cite{mao_artificial_2026}.



\subsection{Comparación empírica: modelos de descriptores frente a GNNs}
\label{ssec:arte_comparacion}


\subsubsection{El debate y sus problemas metodológicos}
 
La pregunta de si las GNNs superan sistemáticamente a los modelos
de descriptores clásicos es una de las más debatidas en quimioinformática moderna. La mayoría de los artículos que proponen nuevas arquitecturas de GNN declaran mejoras sobre la línea base de ML clásico, pero estas comparaciones adolecen frecuentemente de defectos metodológicos: uso de partición aleatoria de una sola corrida, comparación contra versiones sub-optimizadas de los modelos de referencia, o selección de conjuntos de evaluación favorables para la GNN propuesta.
 
Jiang, Wu, Hsieh y colaboradores (2021) publicaron en\textit{Journal of Cheminformatics} el estudio comparativo másriguroso disponible hasta la fecha: ocho algoritmos evaluados sobre11 conjuntos estándar de MoleculeNet con un protocolo de 50 corridasindependientes (cada una con una semilla de partición distinta) paraestimar el rendimiento con su varianza real~\cite{jiang_could_2021}.
 
\subsubsection{Protocolo y conjuntos evaluados}
 
Los cuatro modelos de descriptores (SVM, XGBoost, RF y DNN) emplearon una combinación de 206 descriptores MOE 1D/2D y 1{,}188 huellas (PubChem + subestructura), con optimización de hiperparámetros porTPE (\textit{Tree of Parzen Estimators}) en 50 evaluaciones yselección del mejor modelo sobre el conjunto de validación. Los cuatromodelos de grafos (GCN, GAT, MPNN y Attentive~FP) utilizaron vectores de características atómicas y de enlace estándar, con la misma estrategia de optimización de hiperparámetros.
 
Los autores desarrollaron un script basado en MOE y RDKit para eliminar sales, inorgánicos, contraiones, disolventes y duplicados con etiquetas inconsistentes. Este paso reveló que algunos conjuntos estándar de MoleculeNet contienen hasta un 4\% de estructuras incorrectas o inapropiadas, lo que subraya la necesidad de un preprocesamiento riguroso antes de cualquier modelización~\cite{jiang_could_2021}.
 
\subsubsection{Resultados: cuándo ganan los descriptores y cuándo las GNNs}
 
El hallazgo central puede formularse de forma precisa: \emph{los modelos de descriptores ocupan el 73\% de las posiciones top-3 en los 11 conjuntos evaluados} (24 de 33 posiciones). En términos de rendimiento absoluto, SVM logra las mejores predicciones en dos de los tres conjuntos de regresión (ESOL: RMSE $= 0.569$; FreeSolv: RMSE $= 0.852$), mientras que Attentive~FP lidera en Lipop (RMSE $= 0.553$)~\cite{jiang_could_2021}.
 
Para las tres tareas de clasificación de tarea única, los resultados son inequívocos: en BBBP (2,035 compuestos, clasificación de penetración BHE), RF alcanza AUC-ROC $= 0.927 \pm 0.025$ en el conjunto de prueba promediado sobre 50 corridas, frente a AUC-ROC $= 0.887 \pm 0.032$ de Attentive~FP. En BACE (1,513 compuestos), el patrón se repite: SVM obtiene $0.893 \pm 0.020$, mientras que la mejor GNN (GCN) alcanza $0.898 \pm 0.019$ una diferencia estadísticamente marginal~\cite{jiang_could_2021}.
 
La situación se invierte en los conjuntos grandes y multitarea. En Tox21 (7,811 compuestos, 12 tareas de toxicidad simultáneas), Attentive~FP lidera con AUC-ROC $= 0.852 \pm 0.012$, seguido de RF con $0.838 \pm 0.011$. En ToxCast (8\,539 compuestos, 182 tareas), Attentive~FP alcanza AUC-ROC $= 0.794 \pm 0.017$ frente a RF con $0.782 \pm 0.005$. El mecanismo es el aprendizaje multitarea: la capacidad de las GNNs de compartir representaciones entre endpoints correlacionados les permite generalizar mejor que los modelos univariados de descriptores cuando los datos son abundantes~\cite{jiang_could_2021}.
 
\subsubsection{Coste computacional e interpretabilidad}
 
La disparidad en eficiencia computacional es sustancial. En el conjunto BBBP, XGBoost entrena en $0.242$ y RF en $0.873$; Attentive~FP requiere $98.743$ con GPU (NVIDIA RTX 2080 Ti), y MPNN llega a $316$. En conjuntos grandes como HIV (40,748 compuestos), la diferencia se amplía: SVM tarda $852$, mientras que las GNNs requieren entre $677$ (Attentive~FP) y $1867$ (MPNN)~\cite{jiang_could_2021}. Estas cifras corresponden a hardware de 2020 y deben actualizarse, pero el orden de magnitud de la diferencia es robusto.
 
En cuanto a la interpretabilidad, los autores demostraron que los descriptores relevantes identificados por SHAP sobre XGBoost en ESOL y BBBP son consistentes con el conocimiento fisicoquímico establecido: el logS de Delaney como predictor positivo de solubilidad, la PSA y el número de donores de enlace de hidrógeno como predictores negativos de permeación BHE. Esta capacidad de recuperar conocimiento establecido proporciona una forma de validación mecanística que las GNNs cuyas representaciones atómicas aprendidas no son directamente interpretables no ofrecen de forma inmediata~\cite{jiang_could_2021}.



\subsection{Predicción de permeabilidad a la barrera hematoencefálica}
\label{ssec:arte_bhe}


\subsubsection{La BHE como cuello de botella para el desarrollo de fármacos del SNC}
 
La barrera hematoencefálica es una estructura multicelular formada por células endoteliales de los capilares cerebrales con unionesestrechas, pericitos y los procesos de los astrocitos. Su función fisiológica es proteger el cerebro de sustancias potencialmente tóxicas y mantener la homeostasis del microentorno neuronal, para lo cual restringe el pasaje
paracelular de la mayor parte de las moléculas y exige que los fármacos entren por difusión pasiva o mediante transportadores específicos.
 
Esta restricción tiene consecuencias directas para el descubrimiento de fármacos: el fracaso de candidatos prometedores en ensayos clínicos del SNC se debe en una proporción relevante a permeabilidad BHE insuficiente detectada tardíamente~\cite{spielvogel_enhancing_2025a}. La evaluación experimental estándar log$BB$ in vivo en roedor es costosa, lenta y sujeta a variabilidad interlaboratorio. Los métodos computacionales han sido desarrollados precisamente para anticipar esta propiedad en las fases tempranas del proceso de optimización.
 
\subsubsection{Evolución histórica de los modelos}
 
El trabajo de Faramarzi, Kim, Volpe y colaboradores (2022), desde la FDA, ofrece la síntesis histórica más completa de esta evolución. Su Tabla~1 documenta 30 años de modelos de predicción de log$BB$, comenzando con Young et al. (1988), que ajustaron un modelo de regresión lineal múltiple (MLR) sobre 20 compuestos con $R^2 = 0.69$,hasta aproximaciones modernas de \textit{machine learning} sobre cientos de moléculas~\cite{faramarzi_development_2022}.
 
La primera generación (1988–2005) empleó MLR y mínimos cuadrados parciales (PLS) sobre conjuntos de 20 a 88 compuestos. Estos modelos identificaron los principales determinantes de la permeación BHE pasiva —lipofilia (log$P$), área de superficie polar topológica (TPSA), peso molecular, capacidad de formar enlaces de hidrógeno y establecieron las reglas heurísticas que siguen siendo referencias clínicas: la tPSA menor de 90$^2$ y la ausencia de más de tres pares donores-aceptores como criterios de inclusión para candidatos a SNC~\cite{spielvogel_enhancing_2025a}.
 
La segunda generación (2005–2015) amplió los conjuntos a cientos de compuestos e incorporó métodos no lineales: redes neuronales artificiales (ANN), máquinas de soporte vectorial (SVM) y métodos bayesianos. Esta evolución mejoró la capacidad de capturar interacciones no lineales entre descriptores, aunque los conjuntos siguieron siendo relativamente pequeños y la mezcla de datos \textit{in vivo} e \textit{in vitro} sin discriminación comprometió la homogeneidad de los conjuntos de entrenamiento.
 
\subsubsection{El trabajo de Faramarzi et al. (2022): modelos QSAR
  de referencia de la FDA}
 
El estudio de Faramarzi y colaboradores (2022) representa el límite del estado del arte en modelos estadísticos clásicos para log$BB$. Su conjunto de entrenamiento integra 921 compuestos con datos de permeabilidad BHE \textit{in vivo} obtenidos de literatura publicada, paquetes de aprobación de fármacos no propietarios de la FDA y la base Drug Interaction Database de la Universidad de Washington la colección más amplia de datos log$BB$ in vivo reportada hasta esa fecha \cite{faramarzi_development_2022}.
 
Los autores desarrollaron dos modelos estadísticos independientes sobre plataformas distintas (MultiCASE y Derek), con el objetivo de aprovechar la diversidad metodológica para mejorar la cobertura combinada. Los resultados de validación cruzada fueron sensibilidad de 82–85\% y predictividad negativa (NPV) de 80–83\%; la validación externa sobre 83 compuestos adicionales produjo sensibilidades de 70–75\%, NPV de 70–72\% y cobertura de 78–86\% por modelo individual. Combinando las predicciones de ambos modelos, la cobertura aumentóal 93\%, con la lógica de que los compuestos no cubiertos por un modelo son capturados por el otro~\cite{faramarzi_development_2022}.
 
Desde el punto de vista del dominio de aplicabilidad, los modelos de la FDA tienen una fortaleza importante: están diseñados para fármacos reales (no simplemente moléculas de investigación), lo que los hace relevantes para proyectos orientados a candidatos con perfil farmacológico realista.
 
\subsubsection{El trabajo de Spielvogel et al. (2025): ML multiparamétrico con XAI}
 
Spielvogel, Schindler, Schröder y colaboradores (2025) representan la tercera generación de modelos: ML multiparamétrico sobre una base de datos estandarizada y homogénea, con metodología de validación robusta y análisis de importancia de variables ~\cite{spielvogel_enhancing_2025a}.
 
La base de datos que construyeron es inusualmente rigurosa: 154 moléculas radiolabeladas y fármacos licenciados caracterizados en el mismo laboratorio, con datos experimentales de permeación BHE determinados principalmente por PET (\textit{Positron Emission Tomography}) en modelos humanos y preclínicos la fuente de \textit{ground truth} más directa disponible para permeación CNS in vivo. Las 24 variables de entrada combinan propiedades calculadas (tPSA, log$P$, pKa, peso molecular, PSA 3D una novedad metodológica introducida por los propios autores mediante DFT con funcional B3LYP y base 6-31G(d) con propiedades experimentales (log$P$ cromatográfico HPLC, coeficiente de membrana artificial IAM, unión a albúmina sérica).
 
La validación empleó Monte Carlo con 100 pliegues de partición aleatoria (80:20) para estimar la distribución del AUC con intervalos de confianza del 95\%. El modelo de bosque aleatoriov alcanzó AUC $= 0.88$ (IC 95\%: 0.87–0.90) en la clasificación binaria CNS+/CNS--, y AUC $= 0.82$ (IC 95\%: 0.81–0.82) en la clasificación de tres clases (penetrante, no penetrante, sustrato de transportador de eflujo). Los análisis SHAP revelaron que las tres variables de mayor peso son el BBB score, la PSA (ACD) y la tPSA  propiedades univariadas que ya se conocían como predictores de BHE, pero cuya combinación no lineal por el RF supera sustancialmente su uso individual. Esto contrasta con las puntuaciones de predicción publicadas previamente: CNS MPO (AUC $= 0.53$), CNS MPO PET (AUC $= 0.51$) y BBB score
(AUC $= 0.68$) fueron superadas de forma inequívoca~\cite{spielvogel_enhancing_2025a}.
 
La limitación principal del estudio es el tamaño reducido del conjunto ($n = 154$): aunque la validación Monte Carlo de 100 pliegues controla la varianza estimada, el intervalo de confianza del AUC refleja incertidumbre real. Adicionalmente, la base está dominada por radiofármacos PET, que constituyen una clase estructural específica que puede no representar la diversidad de candidatos orales para el SNC.
 
\subsection{Predicción de toxicidad ADMET multitarea}
\label{ssec:arte_toxicidad}



\subsubsection{La toxicidad como causa dominante de fracaso en desarrollo clínico}
 
La evaluación de toxicidad ocupa una posición central en el \textit{pipeline} de desarrollo farmacéutico: aproximadamente el 30\% de los candidatos preclínicos fracasan por problemas de toxicidad no anticipados, y cerca del 30\% de los fármacos retirados del mercado después de su aprobación lo hacen por toxicidad inesperada~\cite{zhang_computational_2025}. La toxicidad es un concepto multidimensional que abarca la toxicidad aguda (LD$_{50}$), la toxicidad orgánica específica (hepatotoxicidad, cardiotoxicidad, nefrotoxicidad), la genotoxicidad, la carcinogenicidad y la fototoxicidad, entre otros endpoints. Ningún ensayo experimental aislado cubre este espectro, y los ensayos animales tienen costesde tiempo y recursos que los hacen incompatibles con el ritmo de la química medicinal moderna. La toxicidad computacional emerge como respuesta a esta presión: modelos entrenados sobre datos de ensayos de alto rendimiento (\textit{high-throughput screening}) permiten anticipar el riesgo tóxico de miles de candidatos en segundos, con un coste marginal prácticamente nulo una vez construido el modelo.
 
\subsubsection{De las plataformas ADMET a los modelos multitarea}
 
Zhang, Li, Zhang y colaboradores (2025) revisaron más de 20 plataformas de predicción ADMET publicadas hasta 2025, identificando tres generaciones metodológicas~\cite{zhang_computational_2025}:
 
La \textbf{primera generación} corresponde a modelos basados en reglas fisicoquímicas (regla de Lipinski, filtros de Veber, PAINS) y a modelos estadísticos de un solo endpoint. Estas herramientas ofrecen alta velocidad e interpretabilidad, pero su poder predictivos limitado para compuestos que violan las reglas por razones estructurales no contempladas en ellas.
 
La \textbf{segunda generación} empleó ML clásico SVM, RF, redesneuronales superficiales sobre descriptores fisicoquímicos o huellas moleculares. Las iniciativas de referencia son el reto Tox21 (2014) y la plataforma DeepTox, que demostró las ventajas del aprendizaje profundo sobre ML clásico en la predicción multitarea de toxicidad para 12 endpoints simultáneos.
 
La \textbf{tercera generación} aplica GNNs con aprendizaje multitarea, modelos multimodales que integran estructura química con datos de transcriptómica o fenómica, y modelos de lenguaje (\textit{Large
Language Models}) para minería de literatura y predicción
molecular~\cite{zhang_computational_2025}. La tendencia dominante es la transición de la predicción de un único endpoint a la modelización conjunta que explota las correlaciones entre endpoints de toxicidad: la hepatotoxicidad y la cardiotoxicidad no son independientes desde el punto de vista mecanístico, y los modelos que aprovechan esta correlación generalizan mejor~\cite{zhang_computational_2025}.
 
\subsubsection{El estado actual en cardiotoxicidad y toxicidad
  hepática}
 
Los avances más documentados corresponden a los endpoints con mayor disponibilidad de datos. Para la cardiotoxicidad mediada por inhibición del canal hERG uno de los principales predictores de síndrome de QT largo y arritmias fatales , Mao y Yan (2026) reportan que los mejores modelos actuales, basados en GNNs con huellas moleculares, alcanzan AUC-ROC superior a 0.90 en conjuntos de prueba independientes sobre datos de ChEMBL y PubChem~\cite{mao_artificial_2026}. La integración de estos modelos en las
cascadas de selección de candidatos de las compañías farmacéuticas ha permitido excluir el 30–40\% de los compuestos de alto riesgo en etapas previas a la síntesis, reduciendo a la mitad los experimentos de electrofisiología de parche~\cite{mao_artificial_2026}.
 
Para la toxicidad hepática inducida por fármacos (DILI), modelos multitarea que combinan datos transcriptómicos de hepatocitos de rata con datos humanos logran mayor precisión que los modelos entrenados en una sola especie, demostrando el valor de los datos multimodales para la predicción de toxicidad sistémica~\cite{zhang_computational_2025}.
 
\subsubsection{Desafíos metodológicos persistentes}
 
A pesar del progreso, Zhang et al. (2025) identifican tres desafíos que permanecen sin resolver: (\emph{i}) la calidad de los datos de toxicidad es heterogénea mediciones procedentes de distintos laboratorios, protocolos y especies sin anotación sistemática; (\emph{ii}) el desbalance de clases es severo en prácticamente todos los conjuntos de toxicidad, lo que implica
que el AUC-ROC solo puede resultar engañoso y debe complementarse con AUPRC y análisis de matriz de confusión; y (\emph{iii}) los modelos de ML para toxicidad siguen siendo correlacionales, no causales: predicen el riesgo de un señal de toxicidad a partir de similitudes estructurales, no a partir de la comprensión del mecanismo~\cite{zhang_computational_2025}.
 
Esta última limitación es epistemológicamente fundamental y no resuelta por ningún método actual. La distinción entre correlación y causalidad es especialmente relevante cuando los modelos se usan para tomar decisiones de inversión en el desarrollo de nuevas moléculas.



\begin{landscape}

\definecolor{rowA}{HTML}{EBF5FB}
\definecolor{rowB}{HTML}{EAFAF1}
\definecolor{rowC}{HTML}{FEF9E7}
\definecolor{rowD}{HTML}{FDEDEC}
\definecolor{rowE}{HTML}{F4ECF7}
\definecolor{rowF}{HTML}{FDF2E9}

\newcolumntype{P}[1]{>{\RaggedRight\arraybackslash}p{#1}}

\setlength{\tabcolsep}{3pt}
\renewcommand{\arraystretch}{1.18}

\begin{table}[h!]
\centering
\scriptsize

\caption{Síntesis del estado del arte. Seis ejes temáticos con sus
  hallazgos centrales, resultados de referencia y limitaciones
  principales.}
\label{tab:estado_arte}

\begin{tabular}{
  P{2.4cm}
  P{4.2cm}
  P{5.0cm}
  P{3.2cm}
  P{4.6cm}
}

\toprule
\textbf{Área / Referencias} &
\textbf{Enfoque clave} &
\textbf{Hallazgo central} &
\textbf{Resultado} &
\textbf{Limitación principal} \\
\midrule

%--- A ---
\rowcolor{rowA}
\textbf{Afinidad 5-HT1A}\newline
\emph{Czub et al.}\ 2021;\
2022 &
AutoML (DL\,+\,GBM\,+\,XGBoost) sobre 9\,440 ligandos curados
ZINC\,+\,ChEMBL; 216 descriptores Mordred; deduplicación por
$\Delta$pKi. &
Primera base pública diversa con pKi para 5-HT1A. Validación
externa ($>$700 mol., base GLASS) y conformidad OECD. &
$R^2 = 0.74$\newline
RMSE $= 0.54$\newline
(CV-10, split aleatorio) &
Partición aleatoria: métricas optimistas. Ki/IC$_{50}$/EC$_{50}$
mezclados. Sin andamio Bemis--Murcko. Sin dominio de
aplicabilidad. \\

\midrule

%--- B ---
\rowcolor{rowB}
\textbf{GNN afinidad droga--proteína}\newline
\emph{Zhang et al.}\ 2023;\newline
\emph{Hosseini et al.}\ 2022 &
SS-GNN: grafo unificado proteína--ligando (umbral distancia),
0.5\,M parámetros, 0.2\,ms/complejo.\newline
DSD\,/\,FiLMv2: GNN surrogato de docking sobre 128\,M moléculas. &
SS-GNN supera SOTA en PDBbind\,v2016 con modelo 10$\times$ menor.
DSD reemplaza docking clásico con $<$3\% error de recuperación. &
$R_p = 0.853$\newline
(+5.2\% sobre SOTA)\newline
DSD: $9.5\times$ speedup\newline
recall $>97\%$ &
SS-GNN requiere estructura 3D del complejo: inviable en cribado de
novo. DSD evaluado solo en receptor $D_4$. \\

\midrule

%--- C ---
\rowcolor{rowF}
\textbf{Descriptores vs.\ GNN}\newline
\emph{Jiang et al.}\ 2021 &
50 corridas con semillas distintas; 8 algoritmos (SVM, XGBoost, RF,
DNN vs.\ GCN, GAT, MPNN, Attentive\,FP) sobre 11 datasets
MoleculeNet. &
Descriptores ganan 73\,\% de posiciones top-3. GNNs superiores
solo en datasets grandes y multitarea (Tox21, ToxCast). En BBBP
(BHE), RF supera a Attentive\,FP. &
BBBP: RF\,$=0.927$\newline
\phantom{BBBP: }AttFP\,$=0.887$\newline
Tox21: AttFP\,$=0.852$\newline
\phantom{Tox21: }RF\,$=0.838$ &
Partición aleatoria en todos los conjuntos. Sin separación por
andamio. Ningún dataset específico de 5-HT1A. \\

\midrule

%--- D ---
\rowcolor{rowC}
\textbf{Permeabilidad BHE}\newline
\emph{Faramarzi et al.}\ 2022;\newline
\emph{Spielvogel et al.}\ 2025 &
Faramarzi (FDA): 2 modelos QSAR, 921 compuestos, logBB in
vivo.\newline
Spielvogel: RF\,+\,24 parámetros (incl.\ PSA-3D experimental),
MC-CV 100 pliegues, SHAP. &
Spielvogel supera CNS\,MPO (0.53) y BBB\,score (0.68). Introduce
PSA-3D como predictor novedoso. Faramarzi: cobertura combinada
93\,\% con dos modelos complementarios. &
AUC $=0.88$ (binario)\newline
AUC $=0.82$ (3 clases)\newline
Sensibilidad FDA:\newline
82--85\,\% (CV) &
Spielvogel: $n=154$, dominado por radiofármacos PET. Ninguno
declara dominio de aplicabilidad. AUPRC no reportada. \\

\midrule

%--- E ---
\rowcolor{rowD}
\textbf{Toxicidad ADMET multitarea}\newline
\emph{Zhang et al.}\ 2025;\newline
\emph{Mao \& Yan}\ 2026 &
GNN multitarea sobre Tox21/ToxCast. Modelos multimodales
(estructura\,+\,transcriptómica) para DILI. Revisión de $>$20
plataformas ADMET. &
GNNs multitarea superiores en $>$5\,000 compuestos explotando
correlaciones entre endpoints. Modelos hERG excluyen 30--40\,\% de
candidatos de alto riesgo antes de síntesis. &
hERG\,(GNN): AUC\,$>0.90$\newline
Tox21\,(AttFP): 0.852\newline
ToxCast\,(AttFP): 0.794 &
Desbalance severo: AUC-ROC sola engañosa; AUPRC necesaria.
Partición aleatoria en benchmarks. Sin aplicación directa a
ligandos 5-HT1A. \\

\midrule

%--- F ---
\rowcolor{rowE}
\textbf{Reproducibilidad y validez}\newline
\emph{Hasselgren \& Oprea}\ 2024;\newline
\emph{Mao \& Yan}\ 2026 &
Crisis de reproducibilidad dual (experimentos\,+\,modelos ML).
Czub\,(2022): único ejemplo de conformidad OECD explícita en el
corpus. Compuestos AI-diseñados en Fase\,I/II. &
Mayoría de modelos sin dominio de aplicabilidad y con partición
aleatoria. IA como potenciador de decisión experta, no como
sistema autónomo de descubrimiento. &
Éxito clínico:\newline
1--8\,\% (con o sin IA)\newline
Andamio por scaffold:\newline
minoría publicaciones &
Modelos correlacionales, no causales. Sin estándar para dominio de
aplicabilidad en GNNs. Brecha de talento ML\,+\,farmacología. \\

\bottomrule
\end{tabular}

\vspace{3pt}
\noindent{\scriptsize
Brechas pendientes: Chemprop (Yang et al., 2019), B3DB (Meng et al., 2021),
Tox21 original (Mayr et al., 2016), MoleculeNet (Wu et al., 2018).}

\end{table}

\end{landscape}